% Generated draft: requires mathematical and source review.
\documentclass[11pt]{article}
\usepackage[margin=25mm]{geometry}
\usepackage{fontspec}
\setmainfont{Latin Modern Roman}
\usepackage{amsmath,amssymb,mathtools,unicode-math}
\setmathfont{Latin Modern Math}
\usepackage{enumitem,needspace,xurl,hyperref}
\setlength{\parindent}{0pt}
\setlength{\parskip}{5pt}
\setlength{\emergencystretch}{3em}
\newcommand{\N}{\mathbb{N}}
\newcommand{\Z}{\mathbb{Z}}
\newcommand{\Q}{\mathbb{Q}}
\newcommand{\R}{\mathbb{R}}
\newcommand{\C}{\mathbb{C}}
\newcommand{\F}{\mathbb{F}}
\newcommand{\A}{\mathbb{A}}
\newcommand{\PP}{\mathbb{P}}
\newcommand{\E}{\mathbb{E}}
\newcommand{\eps}{\varepsilon}
\newcommand{\dd}{\,\mathrm{d}}
\begin{document}
\section*{UnsolvedMath ID: 20002872}
\textbf{Source problem code:} catalogue entry
\subsection{Definitions and mathematical statement}
S. Berhanu's source question takes a smooth real hypersurface $M^{2n-1}\subset\C^n$ with a strong maximum principle for continuous CR functions and asks whether local continuous CR functions are open. The intended maximum principle is for modulus: on a connected local domain, a continuous CR function with a weak local maximum of its absolute value must be constant. Complex values themselves have no order. The conclusion must exclude constant functions, which are never open as maps to $\C$; this exclusion is implicit in the natural interpretation. CR equations for continuous functions are understood distributionally in tangential $(0,1)$ directions. The original remark says the hypothesis cannot hold in $\C^2$.

The aid's extension-free criterion fixes a nonconstant continuous CR function $f$ on a connected neighborhood $U\subset M$. Let $p\in D\Subset U$ with $D$ connected open and
\[
 f(p)\notin f(\partial D),\qquad
 \delta=\min_{x\in\partial D}|f(x)-f(p)|>0.
\]
The reciprocal maximum-principle argument gives the explicit sufficient conclusion
\[
 \{a\in\C:|a-f(p)|<\delta/2\}\subset f(D).
\]
Indeed a missing value in that disk would make $1/(f-a)$ have an interior maximum larger than all boundary values, contrary to the hypothesis and nonconstancy. Existence of such trapping domains inside every prescribed neighborhood gives openness at $p$. If openness fails in some neighborhood, every such relatively compact connected $D$ within it must instead have $f(p)\in f(\partial D)$. One trapped domain alone is not silently substituted for an arbitrarily local openness statement.

The aid also rules out several local image barriers for a nonconstant function: a complementary exterior circular disk tangent at the image value, a supporting-line half-plane barrier, or containment in a locally regular real-analytic arc that can be holomorphically straightened. Reciprocal or holomorphic-composition maximum arguments exclude these cases. Singular or irregular one-dimensional image sets are not automatically covered. If $f$ is $C^1$, a non-open point must satisfy $\operatorname{rank}_\R df_p<2$, since rank two gives a local submersion.

The supplied literature assessment attributes to Berhanu (2020) affirmative cases for functions extending holomorphically across the hypersurface, ambient-holomorphic restrictions, a dense open subset of each strong-maximum-principle hypersurface, and real-analytic tube structures. Their exact hypotheses and dense-open quantifiers are absent from the record. The general remaining continuous case is not claimed solved. Any proposed counterexample must respect the same-value boundary-fiber and image-barrier restrictions, and the literal constant-function and modulus ambiguities must first be removed.
\subsection{Short English statement}
Under the intended nonconstant maximum-modulus interpretation, trapped same-value fibers give a local image disk by a reciprocal argument. The aid also records known special cases and barriers that a counterexample must avoid, but does not prove the unrestricted continuous open-mapping statement.
\subsection{Sources}
\begin{itemize}
\item Local UnsolvedMath registry, permanent ID 20002872, source code catalogue entry; created and updated 14 August 2026. Canonical \texttt{aim-several-complex-variables-notes.json}, zero-based \texttt{notes[29]}, tag \texttt{problem}, item 1 attributed to S. Berhanu.
\item AIM ``Emerging applications of complexity for CR mappings,'' 2010 workshop, section ``CR manifolds and mappings,'' with the source preamble declaring manifolds smooth unless otherwise stated, according to the supplied verification: \url{https://aimath.org/WWN/crmappings/crmappings.pdf}. Wording is marked exact. The original remark says the hypothesis cannot hold in $\C^2$; literature field is empty.
\item Research attempt 1: \texttt{partial\_result}, problem status at run \texttt{partially\_solved}, rigor label \texttt{proved}, novelty confidence low. It cites Berhanu's 2020 work without a full title, theorem statement, or authoritative bibliographic link.
\item Ulam AI / UnsolvedMath Contributors, CC BY 4.0. Machine-generated research aid, not peer reviewed; independently verify before citation.
\end{itemize}
\end{document}
